
\documentclass{article}
\usepackage[active,tightpage]{preview}
\usepackage{amsmath}
\usepackage{fontspec}
\usepackage[russian]{babel} % Язык по умолчанию, может быть переопределен
\usepackage{cancel}
\usepackage{ulem}
\usepackage{xcolor}
\usepackage{hyperref}

\setmainfont{Arial}[
    Path=./Fonts/, Extension=.TTF,
    UprightFont=*, BoldFont=*BD,
    ItalicFont=*I, BoldItalicFont=*BI
]
\pagenumbering{gobble}
\begin{document}
\begin{preview}
latex \end{preview}\newpage\begin{preview} \begin{center}
    \vspace*{1cm}
    {\Huge \bfseries Вода: $\text{H}_2\text{O}$}
    \vspace*{0.5cm}
    {\Large \itshape Суть Жизни}
\end{center}

\vspace{1cm}

\section{Введение}
Вода ($\text{H}_2\text{O}$) — ключевое вещество на Земле, покрывающее $71\%$ её поверхности и составляющее основу всех живых организмов. Её уникальные свойства критически важны для жизни и всех природных процессов.

\end{preview}\newpage\begin{preview}

\section{Строение Молекулы}
Молекула воды состоит из одного атома кислорода и двух атомов водорода, соединенных ковалентными связями.

\subsection{Геометрия и Полярность}
\begin{itemize}
    \item Молекула имеет угловое (изогнутое) строение с углом Н--О--Н около $104.5^\circ$.
    \item Кислород более электроотрицателен, что создает частичные заряды: $\delta^-$ на кислороде и $\delta^+$ на водороде. Это делает молекулу \textbf{полярной}.
    \item Полярность приводит к образованию \textbf{водородных связей} между молекулами воды, определяющих её уникальные свойства.
\end{itemize}

\begin{center}
    \begin{tikzpicture}[scale=1]
        % Кислород
        \node[circle, draw=red!70!black, fill=red!20, minimum size=1.5cm, label=center:{\Large O}] (O) at (0,0) {};
        % Водород 1
        \node[circle, draw=blue!70!black, fill=blue!20, minimum size=1cm, label=center:{\large H}] (H1) at (-1.5, -1) {};
        % Водород 2
        \node[circle, draw=blue!70!black, fill=blue!20, minimum size=1cm, label=center:{\large H}] (H2) at (1.5, -1) {};
        % Связи
        \draw[line width=1.5pt, black] (O) -- (H1);
        \draw[line width=1.5pt, black] (O) -- (H2);
        % Угол
        \pic [draw, &lt;-&gt;, angle radius=0.8cm, &quot;$\approx 104.5^\circ$&quot;] {angle = H1--O--H2};
    \end{tikzpicture}
    \captionof{figure}{Угловое строение молекулы воды}
\end{center}

\end{preview}\newpage\begin{preview}

\section{Ключевые Свойства Воды}
Водородные связи придают воде ряд аномальных свойств:
\begin{itemize}
    \item \textbf{Высокие температуры кипения/плавления:} Необычно высокие для низкой молекулярной массы.
    \item \textbf{Высокая теплоемкость:} Способность поглощать много тепла без значительного изменения температуры (важно для климата и терморегуляции).
    \item \textbf{Аномалия плотности:} Максимальная плотность при $4^\circ\text{C}$. Лед менее плотен и плавает, предотвращая полное замерзание водоемов.
    \item \textbf{Универсальный растворитель:} Растворяет множество веществ благодаря своей полярности, обеспечивая реакции в живых организмах.
    \item \textbf{Высокое поверхностное натяжение:} Важно для капиллярных явлений и жизни на поверхности воды.
\end{itemize}

\section{Значение для Жизни и Планеты}
Вода незаменима:
\begin{itemize}
    \item \textbf{Биологические функции:} Основной компонент клеток, среда для метаболизма, транспорт веществ, терморегуляция.
    \item \textbf{Экология и климат:} Регулирует глобальные температуры, участвует в круговороте воды, формирует ландшафты.
\end{itemize}

\section{Заключение}
Вода — фундаментальное вещество, чьи уникальные свойства, обусловленные полярностью и водородными связями, делают её основой жизни на Земле. Понимание и сохранение водных ресурсов жизненно важны.

\vspace{1cm}
\begin{center}
    {\itshape Вода — это не просто ресурс, это сама жизнь.&quot;}
\end{center}
\end{preview}\end{document}